\section{Discussion}
\label{sec:discussion}

\subsection{When to use each model class}

The results in Sections~\ref{sec:apriori}--\ref{sec:cost} suggest a practical decision framework for selecting a turbulence closure:

\begin{itemize}
    \item \textbf{Algebraic / transport models} ($k$-$\omega$ SST, EARSM): For applications where computational budget is the primary constraint---design optimization requiring thousands of evaluations, parametric sweeps, or real-time simulation. These models add less than 10\% overhead on GPU and remain the most robust choice for attached or mildly separated flows.

    \item \textbf{MLP (small)}: A reasonable first step toward data-driven closures when the user wants modest accuracy improvement without dramatic cost increase. At 24--50\% overhead, the MLP is affordable for most applications, though its scalar output limits its ability to capture Reynolds stress anisotropy.

    \item \textbf{PI-TBNN}: The recommended choice when the target geometry may differ from the training set---the common case in engineering practice. Despite 42\% worse validation accuracy than the unconstrained TBNN, the PI-TBNN's test RMSE (0.142) is $2.7\times$ better, with particularly strong robustness on the unseen CBFS geometry (CBFS/PH ratio of $1.55\times$ vs.\ $6.05\times$ for TBNN). The computational cost is identical to the TBNN.

    \item \textbf{TBNN}: Should be used with caution. While it achieves the best validation accuracy among neural networks (RMSE 0.085), it exhibits catastrophic overfitting on unseen geometries (test RMSE 0.389). Only recommended when the deployment geometry closely matches the training data.

    \item \textbf{TBRF}: Not recommended for online use but valuable as an offline reference. The TBRF achieves the best test RMSE (0.125) with good generalization ($1.97\times$ degradation), confirming that the accuracy gap between neural networks and tree ensembles persists on unseen geometries. The ensemble averaging provides natural regularization analogous to the PI-TBNN's penalty.
\end{itemize}

\subsection{PI-TBNN: realizability as regularization}

The realizability penalty in the PI-TBNN tells two very different stories depending on whether one examines validation or test performance (Section~\ref{sec:pi_sweep} and~\ref{sec:generalization}).
On the validation set, the penalty appears harmful: PI-TBNN's RMSE (0.120) is 42\% worse than the unconstrained TBNN (0.085).
On the test set, the conclusion reverses: PI-TBNN (0.142) is $2.7\times$ better than TBNN (0.389).

This dramatic reversal reveals that the realizability penalty acts primarily as a \emph{regularizer}, not as a physics constraint.
The TBNN architecture already embeds significant physical structure:
\begin{itemize}
    \item \textbf{Galilean invariance}: The input features (5 invariants of $\Shat$ and $\Omhat$) are frame-indifferent by construction.
    \item \textbf{Symmetry and zero trace}: The tensor basis contraction $b_{ij} = \sum_n g^{(n)} T^{(n)}_{ij}$ guarantees symmetric, traceless output.
    \item \textbf{Near-realizability}: Only 1.3\% of TBNN validation predictions violate realizability bounds.
\end{itemize}

Given that realizability violations are already rare, the penalty's main effect is to constrain the network's effective capacity.
By penalizing outputs near the boundary of the realizable region, it prevents the coefficient functions $g^{(n)}(\lambda_1,\ldots,\lambda_5)$ from developing sharp, geometry-specific features that fit the training data but do not transfer.
The result is a smoother mapping from invariants to tensor coefficients---higher bias but dramatically lower variance.

This reinterpretation has practical implications.
The earlier conclusion of \citet{Wu2018}---that loss-based physics penalties are redundant given architectural constraints---holds for in-distribution accuracy but is misleading for generalization.
When the deployment domain differs from training (as is typical in engineering practice), the PI-TBNN's regularization effect is valuable precisely because the penalty is ``redundant'' for the training data: it constrains the network without distorting the loss landscape in regions where data is available.

\subsection{TBNN overfitting to training geometries}

The TBNN's catastrophic test-set degradation ($4.6\times$ increase in RMSE) is concentrated on the CBFS geometry, where it achieves an RMSE of 0.457 compared to 0.076 on PH $\alpha=1.2$ (Table~\ref{tab:apriori_percase}).
The PH case interpolates between training geometries ($\alpha=0.5$, $0.8$, $1.0$, $1.5$), while the CBFS introduces a qualitatively different separation mechanism (sharp backward-facing step with streamline curvature).

The component-wise analysis (Table~\ref{tab:component_rmse}) reveals that the overfitting is concentrated in the normal stress components ($b_{11}$, $b_{22}$), which show $5$--$7\times$ degradation from validation to test.
The shear stress $b_{12}$, while also degraded, fares better ($4.5\times$).
This pattern is consistent with the TBNN learning geometry-specific normal stress magnitudes---which vary strongly between flow configurations---rather than the universal relationship between strain/rotation invariants and stress anisotropy.

The TBNN has only 9,354 parameters (identical to the PI-TBNN), so the overfitting is not driven by excessive model capacity in the traditional sense.
Rather, the unconstrained loss function allows the network to exploit correlations between invariant features and geometry-specific stress magnitudes that happen to be present in the training set but do not generalize.
The PI-TBNN's realizability penalty disrupts exactly these spurious correlations.

\subsection{TBRF: offline ceiling, online impractical}

The TBRF achieves the highest offline accuracy (RMSE = 0.064 vs.\ 0.085 for TBNN, a 24.6\% improvement) but represents a fundamentally different computational paradigm:
\begin{itemize}
    \item Decision tree traversals involve \emph{data-dependent branching}: different cells follow different paths through the tree, causing warp divergence on GPU.
    \item Tree node data (feature indices, thresholds, child pointers) is scattered in memory with no spatial locality, causing cache misses.
    \item The 10-tree model requires \SI{56}{MB} of tree data, exceeding GPU L2 cache capacity.
\end{itemize}

These are not implementation limitations that can be overcome with better engineering; they reflect a fundamental mismatch between the random forest algorithm and GPU hardware.
The TBRF's value extends beyond establishing a validation accuracy ceiling.
On the test set, the TBRF achieves the best RMSE (0.125) among all models, with a val-to-test degradation ratio of $1.97\times$---far better than the TBNN ($4.60\times$) though not as robust as the PI-TBNN ($1.18\times$).
The ensemble averaging inherent in random forests provides natural regularization, analogous to the PI-TBNN's realizability penalty but through a different mechanism (model averaging vs.\ capacity constraint).
This suggests that the TBRF's offline accuracy advantage over neural networks is partly an accuracy advantage and partly a generalization advantage---it fits the training data better \emph{and} transfers better to new geometries.

\subsection{CPU vs.\ GPU cost differences}

\todo{Discuss how the cost ranking changes between GPU and CPU:
\begin{itemize}
    \item On CPU, matrix multiplies benefit from optimized BLAS libraries, potentially making MLP/TBNN relatively cheaper.
    \item On CPU, tree traversals do not suffer from warp divergence, potentially making TBRF relatively less expensive.
    \item The GPU/CPU speedup ratio may differ across model types: stencil-based models may show higher speedup than NN models.
\end{itemize}
}

\subsection{Limitations}

Several limitations of this study should be noted:

\begin{itemize}
    \item \textbf{Training data dependence}: All NN models are trained on $k$-$\omega$ SST fields from the McConkey dataset. Performance may differ with other RANS baselines ($k$-$\varepsilon$, Spalart--Allmaras) or different training data sources.

    \item \textbf{Generalization}: The test set includes only two cases (periodic hills at $\alpha=1.2$ and CBFS at $\Rey=13{,}700$). While these reveal significant overfitting (Section~\ref{sec:generalization}), performance on dramatically different geometries (3D bluff bodies, internal flows, free shear layers) is not evaluated. The TBNN's catastrophic CBFS failure may or may not be representative of its behavior on other unseen geometries.

    \item \textbf{Single solver}: All cost measurements are from one solver implementation. Different solvers (spectral element, finite volume, unstructured) may show different relative overheads due to different baseline costs.

    \item \textbf{Inference optimization}: Our NN inference uses explicit matrix-vector products. Optimized inference engines (TensorRT, ONNX Runtime) could potentially reduce NN overhead, though the qualitative ordering is unlikely to change.

    \item \textbf{No online training}: We evaluate only fixed-weight, offline-trained models. Online or adaptive training approaches may offer different cost-accuracy tradeoffs.

    \item \textbf{Single GPU type}: Primary results are on the L40S. The relative cost may differ on other GPU architectures (different L2 cache sizes, different FLOP/bandwidth ratios). \todo{Add results on H100/H200 in supplement.}
\end{itemize}

\subsection{Comparison to prior work}

\todo{Discuss how our a priori accuracy compares to reported values in \citet{Ling2016} and \citet{Kaandorp2020}. Note differences in dataset, split protocol, and architecture. Discuss whether the cost ranking changes in the a posteriori setting vs.\ a priori accuracy ordering.}
